% GENERATED FILE. Edit canonical lessons, then run npm run generate:books.
% canonical-source-hash: fnv1a64:d6c9c584b697dd26
% canonical-lessons: BN-C292-khora, BN-C292-durbhikkho, BN-C292-onahar, BN-C292-phata, BN-C292-jholse-jaoya

\chapter{Drought and Hunger}
\label{ch:bn-drought-and-hunger}
\hlchaptermodality{292}

\begin{chapteropening}
\textbf{By the end of this chapter:} \emph{I can talk about a drought, a famine and dry land.}

Complete the last lesson of chapter 292: I can talk about a drought, a famine and dry land.
\end{chapteropening}

\section[\texorpdfstring{\bn{খরা} (khôrā)}{খরা (khôrā)}]{\bn{খরা} (khôrā) — a drought}

\label{lesson:BN-C292-khora}



\begin{quote}
Before the new one: say the Bengali for to hand out, then the Bengali for to donate.
\end{quote}

\subsection*{You'll want to know: \bn{খরা}}
\textbf{\bn{খরা}} — \emph{khôrā} — \textquotedblleft{}a drought\textquotedblright{}.

\begin{quote}
\bn{এই} \bn{বছর} \bn{খরার} \bn{জন্য} \bn{ধান} \bn{হল} \bn{না।}
\end{quote}

\begin{quote}
\emph{ei bôchhôr khôrār jônyô dhān hôlô nā.}
\end{quote}

\begin{quote}
This year the rice failed because of the drought.
\end{quote}

\begin{cousinweb}
\textbf{\bn{খরা}} (\emph{khôrā}) is filed in CDIAL 3819 under Sanskrit \emph{khara}, 'hard, sharp', also 'hot' of a wind, which lists Bengali \emph{kharā}, 'hot and dry', and as a verb 'to parch'. The noun 'a drought', the parching season, is the Bengali use.
\end{cousinweb}

\begin{grammarlens}[title={What you've built}]
One more word for telling what happened.
\end{grammarlens}

\subsection*{Your turn}
\begin{itemize}
\raggedright
  \item \emph{Say it:} \emph{khôrā}
  \item \emph{Say it:} \emph{khôrā}, once more
  \item \emph{Say it:} \emph{ei bôchhôr khôrār jônyô dhān hôlô nā.} — This year the rice failed because of the drought.
  \item \emph{Recall:} say the Bengali for to hand out, then the Bengali for to donate, then say \emph{khôrā} again
\end{itemize}

\begin{tcolorbox}[breakable,colback=teal!4,colframe=teal!35!black,title={Before you move on}]
What does \bn{খরা} mean? (\textquotedblleft{}A drought\textquotedblright{}.) Say it once more.
\end{tcolorbox}
\section[\texorpdfstring{\bn{দুর্ভিক্ষ} (durbhikkhô)}{দুর্ভিক্ষ (durbhikkhô)}]{\bn{দুর্ভিক্ষ} (durbhikkhô) — a famine}

\label{lesson:BN-C292-durbhikkho}



\begin{quote}
Before the new one: say the Bengali for to donate, then the Bengali for a drought.
\end{quote}

\subsection*{You'll want to know: \bn{দুর্ভিক্ষ}}
\textbf{\bn{দুর্ভিক্ষ}} — \emph{durbhikkhô} — \textquotedblleft{}a famine\textquotedblright{}.

\begin{quote}
\bn{দাদু} \bn{একটা} \bn{দুর্ভিক্ষের} \bn{গল্প} \bn{বললেন।}
\end{quote}

\begin{quote}
\emph{dādu ekṭā durbhikkher gôlpo bôllen.}
\end{quote}

\begin{quote}
Grandfather told a story about a famine.
\end{quote}

\begin{cousinweb}
\textbf{\bn{দুর্ভিক্ষ}} (\emph{durbhikkhô}) is Sanskrit \emph{durbhikṣa}, 'dearth, famine' (Monier-Williams): \emph{dur-}, 'bad', and \emph{bhikṣā}, 'alms', literally a bad time for alms.
\end{cousinweb}

\begin{grammarlens}[title={What you've built}]
One more word for telling what happened.
\end{grammarlens}

\subsection*{Your turn}
\begin{itemize}
\raggedright
  \item \emph{Say it:} \emph{durbhikkhô}
  \item \emph{Say it:} \emph{durbhikkhô}, once more
  \item \emph{Say it:} \emph{dādu ekṭā durbhikkher gôlpo bôllen.} — Grandfather told a story about a famine.
  \item \emph{Recall:} say the Bengali for to donate, then the Bengali for a drought, then say \emph{durbhikkhô} again
\end{itemize}

\begin{tcolorbox}[breakable,colback=teal!4,colframe=teal!35!black,title={Before you move on}]
What does \bn{দুর্ভিক্ষ} mean? (\textquotedblleft{}A famine\textquotedblright{}.) Say it once more.
\end{tcolorbox}
\section[\texorpdfstring{\bn{অনাহার} (ônāhār)}{অনাহার (ônāhār)}]{\bn{অনাহার} (ônāhār) — starvation; going without food}

\label{lesson:BN-C292-onahar}



\begin{quote}
Before the new one: say the Bengali for a drought, then the Bengali for a famine.
\end{quote}

\subsection*{You'll want to know: \bn{অনাহার}}
\textbf{\bn{অনাহার}} — \emph{ônāhār} — \textquotedblleft{}starvation; going without food\textquotedblright{}.

\begin{quote}
\bn{দুর্ভিক্ষে} \bn{অনেক} \bn{মানুষ} \bn{অনাহারে} \bn{মরল।}
\end{quote}

\begin{quote}
\emph{durbhikkhe ônek mānush ônāhāre môrlô.}
\end{quote}

\begin{quote}
Many people died of starvation in the famine.
\end{quote}

\begin{cousinweb}
\textbf{\bn{অনাহার}} (\emph{ônāhār}) is Sanskrit \emph{anāhāra}, 'not taking food' (Monier-Williams): \emph{an-}, 'not', and \emph{āhāra}, 'food, eating'.
\end{cousinweb}

\begin{grammarlens}[title={What you've built}]
One more word for telling what happened.
\end{grammarlens}

\subsection*{Your turn}
\begin{itemize}
\raggedright
  \item \emph{Say it:} \emph{ônāhār}
  \item \emph{Say it:} \emph{ônāhār}, once more
  \item \emph{Say it:} \emph{durbhikkhe ônek mānush ônāhāre môrlô.} — Many people died of starvation in the famine.
  \item \emph{Recall:} say the Bengali for a drought, then the Bengali for a famine, then say \emph{ônāhār} again
\end{itemize}

\begin{tcolorbox}[breakable,colback=teal!4,colframe=teal!35!black,title={Before you move on}]
What does \bn{অনাহার} mean? (\textquotedblleft{}Starvation; going without food\textquotedblright{}.) Say it once more.
\end{tcolorbox}
\section[\texorpdfstring{\bn{ফাটা} (phāṭā)}{ফাটা (phāṭā)}]{\bn{ফাটা} (phāṭā) — to crack; to burst}

\label{lesson:BN-C292-phata}



\begin{quote}
Before the new one: say the Bengali for a famine, then the Bengali for starvation.
\end{quote}

\subsection*{You'll want to know: \bn{ফাটা}}
\textbf{\bn{ফাটা}} — \emph{phāṭā} — \textquotedblleft{}to crack; to burst\textquotedblright{}.

\begin{quote}
\bn{খরায়} \bn{জমির} \bn{মাটি} \bn{ফেটে} \bn{গেল।}
\end{quote}

\begin{quote}
\emph{khôrāy jômir māṭi pheṭe gelô.}
\end{quote}

\begin{quote}
In the drought the soil of the fields cracked.
\end{quote}

Past (he/she …): \bn{সেই} \bn{কলসি} \bn{ফাটল}, \emph{sei kolsi phāṭlô}.

\begin{cousinweb}
\textbf{\bn{ফাটা}} (\emph{phāṭā}) comes down from Sanskrit \emph{sphāṭayati}, 'causes to split' (CDIAL 13825, which lists Bengali \emph{phāṭā}, 'to be split'). The -\bn{এ} form is \bn{ফেটে}: \bn{ফেটে} \bn{গেল}, it cracked.
\end{cousinweb}

\begin{grammarlens}[title={What you've built}]
One more word for telling what happened.
\end{grammarlens}

\subsection*{Your turn}
\begin{itemize}
\raggedright
  \item \emph{Say it:} \emph{phāṭā}
  \item \emph{Say it:} \emph{phāṭā}, once more
  \item \emph{Say it:} \emph{khôrāy jômir māṭi pheṭe gelô.} — In the drought the soil of the fields cracked.
  \item \emph{Recall:} say the Bengali for a famine, then the Bengali for starvation, then say \emph{phāṭā} again
\end{itemize}

\begin{tcolorbox}[breakable,colback=teal!4,colframe=teal!35!black,title={Before you move on}]
What does \bn{ফাটা} mean? (\textquotedblleft{}To crack; to burst\textquotedblright{}.) Say it once more.
\end{tcolorbox}
\section[\texorpdfstring{\bn{ঝলসে} \bn{যাওয়া} (jhôlse jāoyā)}{ঝলসে যাওয়া (jhôlse jāoyā)}]{\bn{ঝলসে} \bn{যাওয়া} (jhôlse jāoyā) — to be scorched; to wither in the heat}

\label{lesson:BN-C292-jholse-jaoya}



\begin{quote}
Before the new one: say the Bengali for starvation, then the Bengali for to crack.
\end{quote}

\subsection*{You'll want to know: \bn{ঝলসে} \bn{যাওয়া}}
\textbf{\bn{ঝলসে} \bn{যাওয়া}} — \emph{jhôlse jāoyā} — \textquotedblleft{}to be scorched; to wither in the heat\textquotedblright{}.

\begin{quote}
\bn{কড়া} \bn{রোদে} \bn{গাছের} \bn{পাতা} \bn{ঝলসে} \bn{গেল।}
\end{quote}

\begin{quote}
\emph{kôṛā rode gāchher pātā jhôlse gelô.}
\end{quote}

\begin{quote}
The leaves of the trees were scorched in the fierce sun.
\end{quote}

Past (he/she …): \bn{সেই} \bn{গাছ} \bn{ঝলসে} \bn{গেল}, \emph{sei gāchh jhôlse gelô}.

\begin{cousinweb}
\textbf{\bn{ঝলসে} \bn{যাওয়া}} (\emph{jhôlse jāoyā}) is the -e form of \bn{ঝলসানো}, to scorch — CDIAL 5352 lists Bengali \emph{jhalsā}, 'scorched', under a Middle Indo-Aryan \emph{jhal-}, 'flash' — before \textbf{\bn{যাওয়া}}.
\end{cousinweb}

\begin{grammarlens}[title={What you've built}]
One more word for telling what happened.
\end{grammarlens}

\subsection*{Your turn}
\begin{itemize}
\raggedright
  \item \emph{Say it:} \emph{jhôlse jāoyā}
  \item \emph{Say it:} \emph{jhôlse jāoyā}, once more
  \item \emph{Say it:} \emph{kôṛā rode gāchher pātā jhôlse gelô.} — The leaves of the trees were scorched in the fierce sun.
  \item \emph{Recall:} say the Bengali for starvation, then the Bengali for to crack, then say \emph{jhôlse jāoyā} again
\end{itemize}

\begin{tcolorbox}[breakable,colback=teal!4,colframe=teal!35!black,title={Before you move on}]
What does \bn{ঝলসে} \bn{যাওয়া} mean? (\textquotedblleft{}To be scorched; to wither in the heat\textquotedblright{}.) Say it once more.
\end{tcolorbox}
